\documentclass[10pt,a4paper]{amsart}
\usepackage[margin=19mm]{geometry}
\usepackage{amsmath,amssymb,mathtools}
\usepackage[scr=rsfs,cal=euler]{mathalfa}
\usepackage{xfrac}
\usepackage[unicode,hidelinks,bookmarks=false]{hyperref}
\newcommand{\ie}{i.e.\ }
\newcommand{\cf}{cf.\ }
\newcommand{\resp}{respectively\ }
\newcommand{\Subsection}{\subsection}
\providecommand{\nobreakdash}{\nobreak\hbox{-}}
\setlength{\emergencystretch}{2em}
\multlinegap0pt
\usepackage{bm}
\usepackage{bbm}
\usepackage{dsfont}
\usepackage{xspace}
\usepackage{cancel}
\usepackage{mathtools}
\usepackage{amsthm}
\makeatletter
\@ifundefined{theorem}{\newtheorem{theorem}{Theorem}}{}
\@ifundefined{definition}{\newtheorem{definition}[theorem]{Definition}}{}
\@ifundefined{lemma}{\newtheorem{lemma}[theorem]{Lemma}}{}
\makeatother
\newcommand{\CH}{\mathsf{CH}}
\newcommand{\forceP}{\mathbb{P}}
\newcommand{\forceQ}{\mathbb{Q}}
\title[The Consistency of the $\bf{\Sigma^1\_3}$-Separation Property]{The Consistency of the $\bf{\Sigma^1\_3}$-Separation Property}
\author{Stefan Hoffelner}
\date{Source: arXiv:1912.11811v6 (CC BY 4.0)}
\newtheorem{corollary}{Corollary}
\allowdisplaybreaks
\providecommand{\ldq}{\textquotedblleft}
\providecommand{\rdq}{\textquotedblright}
\newcommand{\forceR}{\mathbb R}
\newcommand{\ZFP}{\mathsf{ZF}^-}
\begin{document}
\maketitle

{\raggedright \noindent\emph{Source and licence.} The Consistency of the $\bf{\Sigma^1_3}$-Separation Property. Version arXiv:1912.11811v6 (CC BY 4.0), distributed under the Creative Commons Attribution 4.0 International licence, as recorded in the arXiv licence metadata for this exact version and shown on the abstract page \url{https://arxiv.org/abs/1912.11811v6}. Full rights notice: \path{licenses/arxiv-1912.11811-CC-BY-4.0.txt}.\par\smallskip\par}

\noindent\textit{Editorial scope.} Selected lightface theorem: source0.tex 1012--1014; complete author proof: 1016--1146. The prior W/Suslin-tree setup, Code and coding area, local Pi-1-2 reading predicate, mixed-support legality/CH preservation, no-accidental-coding result, formula enumeration and earlier dynamic-legality framework are now retained as author source inputs. Their earlier proofs are not recursively reprinted. All modified n-legal rules, the two iterations, frozen R-n construction, tail-legality/uniqueness/definability lemmas, and the final two definitions of the separator remain inside the complete selected proof; these internal lemmas are not deducted as prerequisites. The reviewed level is H3 for this joint construction. Printed author slips are preserved.\par\smallskip
\par\medskip\noindent\textit{Author context (source0.tex lines 132--138).}\par\smallskip
The notation we use will be mostly standard, we hope. We write $\forceP=(\forceP_{\alpha} \, : \, \alpha < \gamma)$ for a forcing iteration of length $\gamma$ with initial segments $\forceP_{\alpha}$. The $\alpha$-th factor of the iteration will be denoted with $\forceP(\alpha)$. Note here that we drop the dot on $\forceP(\alpha)$, even though $\forceP(\alpha)$ is in fact a $\forceP_{\alpha}$-name of a partial order. 
If $\alpha' < \alpha < \gamma$, then we write $\forceP_{\alpha' \alpha}$ to denote the intermediate forcing of $\forceP$ which happens in the interval $[\alpha',\alpha)$, i.e. $\forceP_{\alpha' \alpha}$ is such that 
$\forceP \cong \forceP_{\alpha'} \ast \forceP_{\alpha' \alpha}$.


We write $\forceP \Vdash \varphi$ whenever every condition in $\forceP$ forces $\varphi$, and make deliberate use of restricting partial orders below conditions, that is, if $p \in \forceP $ is such that $p \Vdash \varphi$, we let $\forceP':= \forceP_{\le p}:=\{ q \in \forceP \, : \, q \le p\}$ and use $\forceP'$ instead of $\forceP$. This is supposed to reduce the notational load of some definitions and arguments.
We also sometimes write $V[\forceP]\models \varphi$ to indicate that for every $\forceP$-generic filter $G$ over $V$, $V[G] \models \varphi$.
\par\medskip\noindent\textit{Author context (source0.tex lines 258--259).}\par\smallskip
We have to first create a suitable ground model $W$ over which the actual iteration will take place. $W$ will be a generic extension of $L$, satisfying $\CH$ and, as stated already earlier, has the property that it contains two $\omega_1$-sequences $\vec{S}=\vec{S^1} \cup \vec{S^2}$ of mutually independent Suslin trees. 
The goal is to add the trees generically, and in a second forcing, use an $L$-definable sequence of stationary subsets of $\omega_1$ to code up the trees. The resulting universe will have the feature that any further outer universe, which preserves stationary subsets, can decode the information written into the $L$-stationary subsets in a $\Sigma_1(\omega_1)$-definable way, and hence has access to the sequence of independent Suslin trees $\vec{S}$. This property can be used to create two $\Sigma_1(\omega_1)$-predicates which are empty first and which can be filled with arbitrary reals $x$, using $\aleph_1$-sized forcings with the countable chain condition. These forcing have the crucial feature that they will be independent of the ground model they live in, a feature we will exploit heavily later on.
\par\medskip\noindent\textit{Author context (source0.tex lines 395--418).}\par\smallskip
We introduced the model $W$ for one specific purpose: the possibility to code up reals into the sequence of definable Suslin trees $\vec{S^1}$ or $\vec{S^2}$ using a method which is not sensitive to its ground model.

For the following, we let $W$ be our ground model, though the definitions will work, and will be used for suitable outer models of $W$ as well. We will encounter this situation as ultimately we will iterate the coding forcings we are about to define.
Let $x \in W$  be an arbitrary real, ;et $m,k \in \omega$, let $(x,m,k)$ denote the real which codes the triple consisting of $x,m$ and $k$ in some fixed recursive way, and let $i \in \{ 1,2\}$. Then we shall define the forcing $\hbox{Code}(x,i)$, which codes the real $x$ into $\aleph_1$-many $\omega$-blocks of $\vec{S^i}$ as a two step iteration:

\[ \hbox{Code}((x,m,k),i):= \mathbb{C} (\omega_1)^L \ast \dot{\mathbb{A}} (\dot{Y}_{x,i}) \]
where the first factor is ordinary $\omega_1$-Cohen forcing, but defined in $L$, and the second factor codes a specific subset of $\omega_1$ denoted with $Y_{(x,m,k),i}$ into a real using almost disjoint coding forcing relative to the canonical, constructible almost disjoint family of reals $D$.
We emphasize, that in iterations of coding forcings, we still fall back to force with $(\mathbb{C} (\omega_1))^L$ as our first factor, that is we never use the $\omega_1$-Cohen forcing of the current universe. Thus, iterating the coding forcings is in fact a hybrid of a product (namely the coordinates where we use $(\mathbb{C}(\omega_1))^L$)  and a finites support iteration (the coordinates where we use the almost disjoint coding forcing). We shall discuss this later in more detail.

We let $g \subset \omega_1$ be a $\mathbb{C} (\omega_1)^L$-generic filter over $W$, and let $\rho: [\omega_1]^{\omega} \rightarrow \omega_1$ be some canonically definable, constructible bijection between
these two sets. We use $\rho$ and $g$ to define the set $h \subset \omega_1$, which eventually shall be the set of indices of $\omega$-blocks of $\vec{S}^i$, where we code up the characteristic function of the real ($(x,m,k)$. Let $h:= \{\rho( g \cap \alpha) \,: \, \alpha < \omega_1 \}$ and let $X \subset \omega_1$ be the $<$-least set  (in some previously fixed well-order of $H(\omega_2)^{W[g]}$ which codes the follwing objects:
\begin{itemize}
\item The $<$-least set of $\omega_1$-branches in $W$ through elments of $\vec{S}$ which code $(x,m,k)$ at $\omega$-blocks which start at values in $h$, that is  we collect $\{ b_{\beta} \subset S_{\beta} \, : \, \beta= \omega \gamma + 2n, \gamma \in h \land n \in \omega \land n \notin (x,m,k) \}$ and  $\{ b_{\beta} \subset S_{\beta} \, : \, \beta= \omega \gamma + 2n+1, \gamma \in h \land n \in \omega \land n \in (x,m,k) \}$.
\item The $<$-least set of $\omega_1 \cdot \omega \cdot \omega_1$-many club subsets through $\vec{R}$, our $\Sigma_1 (\omega_1)$-definable sequence of $L$-stationary subsets of $\omega_1$ from the last section, which are necessary to compute every tree $S_{\beta} \in \vec{S}$ which shows up in the above item, using the $\Sigma_1 (\omega_1)$-formula from the previous section before Lemma 2.10.
\end{itemize}

Note that, when working in $L[X]$ and if $\gamma \in h$ then
 we can read off $(x,m,k)$ via looking at the $\omega$-block of $\vec{S^i}$-trees starting at $\gamma$ and determine which tree has an $\omega_1$-branch in $L[X]$:
\begin{itemize}
 \item[$(\ast)$]  $n \in (x,m,k)$ if and only if $S^i_{\omega \cdot \gamma +2n+1}$ has an $\omega_1$-branch, and $n \notin (x,m,k)$ if and only if $S^i_{\omega \cdot \gamma +2n}$ has an $\omega_1$-branch.
\end{itemize}
Note that $(\ast)$ is actually a formula $(\ast) ((x,y,m) ,\gamma)$ with two parameters $(x,y,m)$ and $\gamma$ but we will suppress this, as the parameters usually are clear from the context.
Indeed if $n \notin (x,m,k)$ then we added a branch through $S^i_{\omega \cdot \gamma+ 2n}$. If on the other hand $S^i_{\omega \cdot\gamma +2n}$ is 
Suslin in $L[X]$ then we must have added an $\omega_1$-branch through $S^i_{\omega \cdot \gamma +2n+1}$ as we always add an $\omega_1$-branch through either $S^i_{\omega \cdot \gamma +2n+1}$ or $S^i_{\omega \cdot \gamma +2n}$ and adding branches through some $S^i_{\alpha}$'s  will not affect that some $S^i_{\beta}$ is Suslin in $L[X]$, as $\vec{S}$ is independent.
\par\medskip\noindent\textit{Author context (source0.tex lines 421--449).}\par\smallskip
We note that we can apply an argument resembling David's trick in this situation. We rewrite the information of $X \subset \omega_1$ as a subset $Y \subset \omega_1$ using the following line of reasoning.
It is clear that any transitive, $\aleph_1$-sized model $M$ of $\ZFP$ which contains $X$ will be able to correctly decode out of $X$ all the information.
Consequentially, if we code the model $(M,\in)$ which contains $X$ as a set $X_M \subset \omega_1$, then for any uncountable $\beta$ such that $L_{\beta}[X_M] \models \ZFP$ and $X_M \in L_{\beta}[X_M]$:
\[L_{\beta}[X_M] \models \text{\ldq The model decoded out of }X_M \text{ satisfies $(\ast)$ for every $\gamma \in h \subset \omega_1$\rdq.} \]
In particular there will be an $\aleph_1$-sized ordinal $\beta$ as above and we can fix a club $C \subset \omega_1$ and a sequence $(M_{\alpha} \, : \, \alpha \in C)$ of countable elementary submodels such that
\[\forall \alpha \in C (M_{\alpha} \prec L_{\beta}[X_M] \land M_{\alpha} \cap \omega_1 = \alpha)\]
Now let the set $Y\subset \omega_1$ code the pair $(C, X_M)$ such that the odd entries of $Y$ should code $X_M$ and if $Y_0:=E(Y)$ where the latter is the set of even entries of $Y$ and $\{c_{\alpha} \, : \, \alpha < \omega_1\}$ is the enumeration of $C$ then
\begin{enumerate}
\item $E(Y) \cap \omega$ codes a well-ordering of type $c_0$.
\item $E(Y) \cap [\omega, c_0) = \emptyset$.
\item For all $\beta$, $E(Y) \cap [c_{\beta}, c_{\beta} + \omega)$ codes a well-ordering of type $c_{\beta+1}$.
\item For all $\beta$, $E(Y) \cap [c_{\beta}+\omega, c_{\beta+1})= \emptyset$.
\end{enumerate}
We obtain
\begin{itemize}
\item[$({\ast}{\ast})$] For any countable transitive model $M$ of $\ZFP$ such that $\omega_1^M=(\omega_1^L)^M$ and $ Y \cap \omega_1^M \in M$, $M$ can construct its version of the universe $L[Y \cap \omega_1^M]$, and the latter will see that there is an $\aleph_1^M$-sized transitive model $N \in L[Y \cap \omega_1^M]$ which models $(\ast)$ for $(x,m,k)$ and every $\gamma \in h \subset \omega_1^M$.
\end{itemize}
Thus we have a local version of the property $(\ast)$.

In the next step $\dot{\mathbb{A}} (\dot{Y})$, working in $W[g]$, for $g\subset \mathbb{C} (\omega_1)$ generic over $W$, we use almost disjoint forcing $\mathbb{A}_D(Y)$ relative to the $<_L$-least almost disjoint family of reals $D \in  L $ to code the set $Y$ into one real $r$. This forcing is well-known, has the ccc and its definition only depends on the subset of $\omega_1$ we code, thus the almost disjoint coding forcing  $\mathbb{A}_D(Y)$ will be independent of the surrounding universe in which we define it, as long as it has the right $\omega_1$ and contains the set $Y$.

We finally obtained a real $r$ such that
\begin{itemize}
\item[$({\ast}{\ast}{\ast})$] For any countable, transitive model $M$ of $\ZFP$ such that $\omega_1^M=(\omega_1^L)^M$ and $ r  \in M$, $M$ can construct its version of $L[r]$ which in turn thinks that there is a transitive $\ZFP$-model $N$ of size $\aleph_1^M$  such that $N$ believes $(\ast)$ for $(x,m,k)$ and every $\gamma \in h$.
\end{itemize}

Note that the above is a $\Pi^1_2(r)$-statement.
We say in this situation that the real $(x,m,k)$\emph{ is written into $\vec{S}^i$}, or that $(x,m,k)$ \emph{is coded into} $\vec{S^i}$.
If $(x,m,k)$ is coded into $\vec{S}^i$ and $r$ is a real witnessing this, then the set $h$ which is equal to $\{ \gamma < \omega_1 \, ;\, \gamma \text{ is a starting point for an $\omega$-block where $(\ast)$ }$ for $(x,y,m)$ holds$\}$ is dubbed (following \cite{FS}) the coding area of $(x,m,k)$ with respect to $r$.
\par\medskip\noindent\textit{Author context (source0.tex lines 451--474).}\par\smallskip
We want to iterate these coding forcings. As the first factor of a coding forcing will always be $(\mathbb{C}(\omega_1))^L$, an iteration of the coding forcing is in fact a hybrid of a (countably supported) product (namely the coordinates where we use $(\mathbb{C}(\omega_1))^L$)  and an actual finite support iteration (the coordinates where we use almost disjoint coding forcing).

\begin{definition}
A mixed support iteration $\forceP=(\forceP_{\beta}\,:\, {\beta< \alpha})$ is called legal if $\alpha < \omega_1$ and there exists a bookkeeping function $F: \alpha \rightarrow H(\omega_2)^2$ such that 
 $\forceP$ is defined inductively using $F$ as follows:
 \begin{itemize}
 \item If $F(0)=(x,i)$, where $x$ is a real, $i\in \{1,2\}$, then $\forceP_0= \hbox{Code}({x,i} )$. Otherwise $\forceP_0$ is the trivial forcing.
 \item If $\beta>0$ and $\forceP_{\beta}$ is defined, $G_{\beta} \subset \forceP_{\beta}$ is a generic filter over $W$, $F(\beta)=(\dot{x}, i)$, where $\dot{x}$ is a $\forceP_{\beta}$-name of a real, $i \in \{1,2\}$ and $\dot{x}^{G_{\beta}}=x$ then, working in $W[G_{\beta}]$ we let
 $\forceP(\beta):= \hbox{Code} ({x,i} )$, that is we code $x$ into the $\vec{S}^i$, using our coding forcing. We shall use full (i.e. countable) support on the $(\mathbb{C}(\omega_1))^L$-coordinates and finite support on the coordinates where we use almost disjoint coding forcing.
 \end{itemize}
 
\end{definition}
 Informally speaking, a legal  forcing just decides to code the reals which the bookkeeping $F$ provides into either $\vec{S^1}$ or $\vec{S^2}$. Note further that the notion of legal  can be defined in exactly the same way over any $W[G]$, where $G$ is a $\forceP$-generic filter over $W$ for an legal  forcing. Finally note that instead of creating $\omega$-blocks of Suslin trees using $\mathbb{C} (\omega_1)^L$ where we code the branches every single time we code a real, we could have also defined an altered ground model $W'$ as $W[g]$, where $g \subset \prod \mathbb{C} (\omega_1)$ is generic for the countably supported product of $\aleph_1$-many copies of $\omega_1$-Cohen forcing, and then worked over $W'$ using exclusively almost disjoint coding forcings which pick first one coordinate $g_{\alpha}$, $\alpha< \omega_1$ of $g$ in an injective way, and then code the $\aleph_1$-many branches along $g_{\alpha}$ using almost disjoint coding forcings as described above. The difference between these approaches is only of symbolic nature, we opted for the one we chose because of a slightly neater presentation. 
 
 We obtain the following first properties of legal  forcings:
 \begin{lemma} 


\begin{enumerate}
\item If $\forceP=(\forceP(\beta) \, : \, \beta < \delta) \in W$ is legal  then for every $\beta < \delta$, $\forceP_{\beta} \Vdash| \forceP(\beta)|= \aleph_1$, thus every factor of $\forceP$ is forced to have size $\aleph_1$.
\item Every legal  forcing over $W$ preserves $\aleph_1$ and $\CH$.
\item The product of two legal  forcings is legal  again.
\end{enumerate}
\end{lemma}
\par\medskip\noindent\textit{Author context (source0.tex lines 484--488).}\par\smallskip
The second assertion of the last lemma immediately gives us the following:
\begin{corollary}
Let $\forceP= (\forceP(\beta) \, : \, \beta < \delta) \in W$ be an legal  forcing over $W$. Then $W[\forceP] \models \CH$. Further, if $\forceP= (\forceP(\alpha) \, : \, \alpha < \omega_1) \in W$ is an $\omega_1$-length iteration such that each initial segment of the iteration is legal  over $W$, then $W[\forceP] \models \CH$.

\end{corollary}
\par\medskip\noindent\textit{Author context (source0.tex lines 490--501).}\par\smallskip
In an iteration of coding forcing we do not add any unwanted or accidental solutions to our $\Sigma^1_3$ predicate give by $({\ast} {\ast} {\ast})$, which we shall show now.

The set of triples of (names of) reals which are enumerated by the bookkeeping function $F \in W$ which comes along with an legal $\forceP = (\forceP(\beta) \, : \, \beta < \delta)$, we call the set of reals coded by $\forceP$. That is, if \[ \forceP(\beta)= (\mathbb{C}(\omega_1))^L \ast \dot{\mathbb{A}}_D (\dot{Y}_{(\dot{x}_{\beta}, \dot{y}_{\beta}, \dot{m}_{\beta} ) } ) \] and $G \subset \forceP$ is a generic filter and if we let for every $\beta < \delta$,
$ \dot{x}_{\beta}^G =:x_{\beta}$, $\dot{y}_{\beta}^G =:y_{\beta}$, $\dot{m}_{\beta}^G =:m_{\beta}$,  then
$\{ (x_{\beta},y_{\beta},m_{\beta} ) \, : \, \beta < \alpha \}$ is the set of reals coded by $\forceP$ and $G$ (though we will suppress the $G$).

\begin{lemma}
If $\forceP \in W$ is legal, $\forceP=(\forceP_{\beta} \, : \, \beta < \delta)$, $G \subset \forceP$ is generic over $W$ and $\{ (x_{\beta},y_{\beta},m_{\beta} ) \, : \, \beta < \delta\}$ is the set of (triples of) reals which is coded as we use $\forceP$. Let 
\[A:= \{ (x,m,k) \in W[G] \, : \, \exists r ( ({\ast}{\ast}{\ast} )\text{ holds for $r$ and $(x,m,k)$} \}.  \] Then
in $W[G]$, the set of reals which belong to $A$ is exactly 
$\{ (x_{\beta},y_{\beta},m_{\beta} ) \, : \, \beta < \delta\}$, that is, we do not code any unwanted information accidentally.
\end{lemma}
\par\medskip\noindent\textit{Author context (source0.tex lines 555--559).}\par\smallskip
For $n \in \omega$ let \[\varphi_n(v_0)= \exists v_1 \psi_n(v_0,v_1)\]  be the $n$-th formula in an enumeration of the $\Sigma^1_3$-formulas with one free variable. Let \[A_n:= \{ x \in 2^{\omega} \, : \, \varphi_n(x) \text{ is true} \},\] so $A_n$ is the set of reals whose definition uses the $n$-th $\Sigma^1_3$-formula in our enumeration. We force with an $\omega_1$-length
mixed support iteration of legal forcings which all have size $\aleph_1$, and use 
a definable, surjective bookkeeping-function \[ F: \omega_1 \rightarrow \omega_1  \times \omega_1 \times \omega \times \omega\]  to determine the iteration. We demand that every $\alpha < \omega_1$ is always strictly bigger than the first projection of $F(\alpha)$. We also assume that every quadruple $(\beta, \gamma, m,k)$ in $\omega_1  \times \omega_1 \times \omega \times \omega$ is hit unboundedly often by $F$.

The purpose of $F$ is to list all triples of the form $(x,m,k)$, where $x$ is a real in some intermediate universe of our iteration and $m,k \in \omega$ corresponds to a pair $(\varphi_m,\varphi_k)$ of $\Sigma^1_3$-formulas. The iteration will be defined in such a way, that, at every stage $\beta$ of the iteration, whenever some triple $(x,m,k)$ is considered by $F$, we must decide immediately whether to code (a real coding) the triple $(x,m,k)$ somewhere into the $\vec{S^1}$ or $\vec{S^2}$-sequence. The set of codes written into $\vec{S}^1$ which contain $m,k$ will result in the $\Sigma^1_3$-set $D^1_{m,k}$ which is a supset of $A_m$, the set of codes containing $m,k$ which are written into $\vec{S^2}$ shall result in the $\Sigma^1_3$-supset $D^2_{m,k}$ of $A_k$. The real $R_0$ will be used to indicate, in a uniform way for all $m,k$, the set of $\aleph_1$-many $\omega$-blocks of $\vec{S}$, which represent insecure data  we should not use for our separating sets.
\par\medskip\noindent\textit{Author context (source0.tex lines 721--757).}\par\smallskip
To motivate the following we first consider a more fine-tuned approach to the definition of the iteration of the proof of the last theorem. Let us assume that $(m,k)$ is the first pair such that case b in the definition of the iteration applies. Recall that in the discussion of case b, we showed that, given a pair $(A_m, A_k)$ of $\Sigma^1_3$-sets for which there does not exist a legal forcing $\forceQ$ such that $\Vdash_{\forceQ} \exists z (\varphi_m(z) \land \varphi_k(z))$ becomes true, we can assign for an arbitrary real $x$ always a side $\vec{S}^1$ or $\vec{S}^2$ such that in all future legal extensions, there will never occur a pathological situation, i.e. from that stage on we never run into the problem of having coded the triple $(x,m,k)$ into, say, $\vec{S^1}$, yet $\varphi_k(x)$ becomes true in some future extension of our iteration (or vice versa). Note here that the arguments in the discussion of case b were uniform for all reals $x$ which appear in a legal extension. So it is reasonable to define for the pair $(m,k)$ 
a stronger notion of legality, called 1-legal with respect to $(0,m,k)$ (the 0 indicates the base case of an inductive construction we define later) as follows:

Let $F: \gamma \rightarrow H(\omega_2)^4$ be a bookkeeping function and let $E:=\{ (0,m,k) \}$. We let $\forceP$ be a mixed support iteration of length $ \gamma$. Then we say that $(\forceP  \,  : \, \beta < \gamma\} ) $ is 1-legal with respect to $E$ and $F$ if
\begin{itemize}
\item $\forceP$ is a legal forcing relative to $F$.
\item Whenever $\beta < \gamma$ is a stage such that $F(\beta)=(\dot{x},m,k,i)$, where $\dot{x}$ is a $\forceP_{\beta}$-name of a real, $\xi$ is an ordinal and $i \in \{1,2 \}$ we split into three subcases:
\begin{itemize}
\item[(i)] First we assume that in $W[G_{\beta}]$,
there is a legal forcing $\forceQ$ such that $\forceQ \Vdash x (=\dot{x}^{G_{\beta}}) \in A_m$.

Then, the $\beta$-th forcing of $\forceP$, $\forceP(\beta)$ must be $\hbox{Code} ((x,m,k), 1)$.
\item[(ii)]
Assume that (i) is wrong but the dual situation is true for $A_k$. That is, 
there is a legal forcing $\forceQ$ such that $\forceQ \Vdash x (=\dot{x}^{G_{\beta}}) \in A_k$.


Then, the $\beta$-th forcing of $\forceP$, $\forceP(\beta)$ must be $\hbox{Code} ((x,m,k), 2)$.
\item[(iii)]

In the third case  we assume that neither (i) nor (ii) is true.
In that situation we force with either coding $x$ into the $A_m$ or the $A_k$ side, whatever the bookkeeping tells us.

%If neither (i) nor (ii) is true, then 


%$\forceP(\beta)=\hbox{Code} ((x,m,k),i,\gamma')$ for $i=F(\beta)_5$ and $\gamma'=F(\beta)_4$. We also demand that $\dot{B}(\beta+1) = \dot{B} (\beta) \cup F(\beta)_3$, provided $F(\beta)_3$ is a $\forceP_{\beta+1}$-name of a countable subset of $\{ 1,2 \} \times \omega_1$.
\end{itemize}

If the iteration  $(\forceP \, : \, \beta < \gamma)$ obeys the above rules, then we say that it is 1-legal with respect to $E$ and bookkeeping function $F$.
\end{itemize}

From earlier considerations it is clear that if we drop the notion of legal from now on and replace it with 1-legal relative to $E=\{ (0,m,k)\}$ in our iteration,  we can ensure that, at least for the pair $(m,k)$ no new pathological situations will arise anymore.
This process can be iterated. Assume that we run into a new pair $(m',k')$ where the modified case b applies, i.e. there is no 1-legal forcing which forces $\exists z (\varphi_{m'}(z) \land\varphi_{k'} (z))$, we can introduce the new notion of 2-legal with respect to $\{(0,m,k) , (1,m',k')\}$. If chosen the right way, this new notion will hand us a condition that guarantees that no new pathological situations arises for the two pairs $(\varphi_m,\varphi_k)$ and $(\varphi_{k'}, \varphi_{m'})$. 

So our strategy for producing a model where the $\bf{\Sigma^1_3}$-separation property is as follows: we list all possible reals $x$, parameters $y$ and pairs of $\Sigma^1_3$-formulas
$(\varphi_m(\cdot,y), \varphi_k(\cdot,y))$, while simultaneously define stronger and stronger versions of legality, which take care of placing the reals we encounter along the iteration in a non-pathological way.
\par\medskip\noindent\textit{Author context (source0.tex lines 767--825).}\par\smallskip
The notions of 0 and 1-legality will form the base cases of an inductive definition.
Let $\alpha \ge 1$ be an ordinal and assume we defined already the notion of $\alpha$-legality. Then we can inductively define the notion of $\alpha+1$-legality  as follows.

Suppose that $\gamma < \omega_1$, $F$ is a bookkeeping function,  \[F: \gamma \rightarrow H(\omega_2)^5 \] and \[\forceP=(\forceP_{\beta} \, : \, \beta < \gamma)\] is a legal forcing relative to $F$ (in fact relative to some bookkeeping $F'$ determined by $F$ in a unique way - the difference here is not relevant). 

Suppose that 
\[E= \{(\delta, \dot{y}_{\delta}, m_{\delta} ,k_{\delta}) \, : \,  \delta \le \alpha\}\]
where $m_{\delta},k_{\delta} \in \omega$ and every $\dot{y}_{\delta}$ is a  $\forceP$-name of a real and for every two ordinals $\beta, \gamma< \alpha$, $\forceP \Vdash (\dot{y}_{\beta}, m_{\beta},k_{\beta}) \ne (\dot{y}_{\gamma},m_{\gamma},k_{\gamma})$. Suppose that for every $\delta \le \alpha$, $(\forceP_{\beta} \, : \, \beta < \gamma )$ is $\delta$-legal with respect to $E \upharpoonright \delta = \{ (\eta, \dot{y}_{\eta},m_{\eta},k_{\eta}) \in E \, : \, \eta < \delta \}$ and $F$.
Finally assume that $\dot{y}_{\alpha+1}$ is a $\forceP$-name for a real and $m_{\alpha+1},k_{\alpha+1} \in \omega$ such that $\forceP \Vdash \forall \delta \le \alpha ((\dot{y}_{\delta},m_{\delta},k_{\delta}) \ne (\dot{y}_{\alpha+1},m_{\alpha+1},k_{\alpha+1}))$. Then we say that $(\forceP_{\beta} \, : \, \beta < \gamma )$ is $\alpha+1$-legal with respect to $E \cup \{\alpha+1, \dot{y}_{\alpha+1},m_{\alpha+1},k_{\alpha+1})\}$ and $F$ if it obeys the following rules. 
\begin{enumerate}
\item Whenever $\beta < \gamma $ is such that there is a $\forceP_{\beta}$-name $\dot{x}$ of a real and an integer $ i\in\{1,2\}$ such that
\[F(\beta)= (\dot{x},\dot{y}_{\alpha+1},m_{\alpha+1},k_{\alpha+1},i)\] and $\dot{y}_{\alpha+1}$ is in fact a $\forceP_{\beta}$-name, and for $G_{\beta}$  a $\forceP_{\beta}$-generic over $W$, $W[G_{\beta}]$ thinks that
\begin{align*}
\exists \forceQ (&\forceQ \text{ is } \alpha \text{-legal with respect to } E \,  \land \\& \forceQ \Vdash x \in A_m({y}_{\alpha+1})),
\end{align*}
where $x=\dot{x}^G$, and $y_{\alpha}=\dot{y}_{\alpha+1}^G$.
Then continuing to argue in $W[G_{\beta}]$, if ${\forceQ}_1=\dot{\forceQ}_1^{G_{\beta}}$ we let
\[\forceP(\beta)= \hbox{Code}((x,y,m,k),1). \]  Note that we confuse here the quadruple $(x,y,m,k)$ with one real $w$ which codes this quadruple.

\item
Whenever $\beta < \gamma $ is such that there is a $\forceP_{\beta}$-name $\dot{x}$ of a real and an integer $ i \in \{1, 2\}$ such that
\[F(\beta)= (\dot{x},\dot{y}_{\alpha+1},m_{\alpha+1},k_{\alpha+1}, i)\] and for $G_{\beta}$ which is $\forceP_{\beta}$-generic over $W$, $W[G_{\beta}]$ thinks that
\begin{align*}
\forall \forceQ_1 (&\forceQ_1 \text{ is } \alpha \text{-legal with respect to } E   \\& \rightarrow \, \lnot(\forceQ_1 \Vdash x \in A_m(\dot{y}_{\alpha+1})))
\end{align*}
but there is a forcing $\forceQ_2$ such that $W[G_{\beta}]$ thinks that 
\begin{align*}
\forceQ_{2} \text{ is } \alpha &\text{-legal with respect to $E$ and } \\ & {\forceQ_2} \Vdash x \in A_k( \dot{y}_{\alpha+1} )
\end{align*}


Then continuing to argue in $W[G_{\beta}]$, we force with
\[\forceP(\beta):= \hbox{Code}((x,y,m,k),2).\]  Note that we confuse here again the quadruple $(x,y,m,k)$ with one real $w$ which codes this quadruple.




\item If neither 1 nor 2 is true, then either \[ \forceP(\beta)=\hbox{Code}((x,y,m,k),2)\] or 
\[ \forceP(\beta)=\hbox{Code}((x,y,m,k),1)\] depending on whether $ i\in\{1,2\}$ in $F(\beta)$ was 1 or 2. 

\item 
If $F(\beta) = (\dot{x},\dot{y},m,k,i)$ and for our $\forceP_{\beta}$-generic filter $G$, $W[G] \models \forall \delta \le \alpha+1 ((\delta,\dot{y}^G,m,k) \notin E^G)$,
then, working over $W[G_{\beta}]$ let
\[ \forceP(\beta)=\hbox{Code}((x,y,m,k),i)\] depending on whether $ i\in\{1,2\}$ in $F(\beta)$ was 1 or 2.
\end{enumerate} 
This ends the definition for the successor step $\alpha \rightarrow \alpha+1$.
For limit ordinals $\alpha$, we say that a legal forcing $\forceP$ is $\alpha$ legal with respect to $E$ and $F$ if for every $\eta < \alpha$,
$(\forceP_{\beta} \,: \, \beta < \gamma)$ is $\eta$-legal with respect to $E \upharpoonright \eta$ and some $F'$.
\par \medskip

We add a couple of remarks concerning the last definition.
\begin{itemize}

\item By definition, if $\delta_2 < \delta_1$ and $\forceP_1$ is $\delta_1$-legal with respect to $E= \{(\beta, \dot{y}_{\beta}, m_{\beta} ,k_{\beta}) \, : \,  \beta \le \delta_1\}$ and some $F_1$, then  $\forceP_1$ is also $\delta_2$-legal with respect to $E \upharpoonright \delta_2 = \{(\beta, \dot{y}_{\beta}, m_{\beta} ,k_{\beta}) \, : \,  \beta \le \delta_2\}$ and an altered bookkeeping function $F'$.
 
\item The notion of $\alpha$-legal can be defined in a uniform way over any legal extension $W'$ of $W$.

\item We will often just say that some iteration $\forceP$ is $\alpha$-legal, by which we mean that there is a set $E$ and a bookkeeping $F$ such that $\forceP $ is $\alpha$-legal with respect to $E$ and $F$.
\end{itemize}
\par\medskip\noindent\textit{Author context (source0.tex lines 829--832).}\par\smallskip
\begin{lemma}\label{productlegal}
Let $\alpha \ge 1$, assume that $W'$ is some $\alpha$-legal generic extension of $W$, and that $\forceP^1=(\forceP^1_{\beta} \,: \,\beta < \delta)$ and $\forceP^2=(\forceP^2_{\beta} \,: \, \beta < \delta) $ are two $\alpha$-legal forcings over $W'$ with respect to a common set $E=\{\delta,\dot{y}_{\delta},m_{\delta}, k_{\delta} \,: \, \delta < \alpha\}$ and bookkeeping functions $F_1$ and $F_2$ respectively. 
Then there is a bookkeeping function $F$ such that $\forceP_1 \times \forceP_2$ is $\alpha$-legal over $W'$ with respect to $E$ and $F$.
\end{lemma}

\begin{theorem}
Starting with $L$ as the ground model, one can produce a set-generic extension $L[G]$ which satisfies $\CH$ and in which the ${\Sigma^1_3}$-separation property holds.
\end{theorem}

\begin{proof}


For the the proof to come, we will redefine the notion of legal forcings.
\begin{definition}
A mixed support iteration is called (0-)legal if it is defined as in Definition 4.2, with the only difference that we code (reals that code) quadruples of the form $(x,0,m,k)$ and $(x,1,m,k)$, where $m,k \in \omega$ and $x$ is (the name of) a real into $\vec{S}$.
\end{definition}
Similar to before, we take 0-legal forcings as the base set of forcings and define gradually smaller families of forcings, which we call $n$-legal forcings with respect to a set $E$ and a bookkeeping function $F$.
Assume that $E$ is a finite list of length $n$ of pairwise distinct pairs of natural numbers $(m_l, k_l)$, $l \le n$,
and that $F$ is a bookkeeping function.
Assume that for every $l \le n$ we do have a notion of $l$-legality with respect to $E \upharpoonright l$, let $(m_{n+1}, k_{n+1})$ be a new pair of natural numbers, distinct from the previous ones. Then we say that $\forceP$ is $n+1$-legal with respect to $E \cup \{ (m_{n+1}, k_{n+1} ) \}$ and $F$ if for every $l \le n$, $\forceP $ is $l$-legal with respect to $E \upharpoonright l$ and $F$ and it obeys the following rules:
\begin{enumerate}
\item $\forceP$ never codes a quadruple of the form $(x,1,m_{n+1},k_{n+1})$ into $\vec{S}$.
\item Whenever $\beta < \gamma $, where $\gamma$ is the length of the iteration $\forceP$, is such that there is a $\forceP_{\beta}$-name $\dot{x}$ of a real and an integer $ i\in\{1,2\}$ such that
\[F(\beta)= (\dot{x},m_{n+1},k_{n+1},i)\] and for  $G$ which is $\forceP_{\beta}$-generic over $W$, $W[G]$ thinks that
\begin{align*}
\exists \forceQ (&\forceQ \text{ is } n \text{-legal with respect to } E \,  \land \\&   \forceQ \Vdash x \in A_{m_{n+1}}),
\end{align*}
where $x=\dot{x}^G$, and $y_{\alpha}=\dot{y}_{\alpha+1}^G$.
Then continuing to argue in $W[G]$, we let
\[\forceP(\beta)= \hbox{Code}((x,0,m_{n+1},k_{n+1}),1).\] 

\item
Whenever $\beta < \gamma $ is such that there is a $\forceP_{\beta}$-name $\dot{x}$ of a real and an integer $ i\in\{1,2\}$ such that
\[F(\beta)= (\dot{x},m_{n+1},k_{n+1},i)\] and for $G_{\beta}$ which is $\forceP_{\beta}$-generic over $W$, $W[G_{\beta}]$ thinks that
\begin{align*}
\forall \forceQ_1 (&\forceQ_1 \text{ is } n \text{-legal with respect to } E  \\& \rightarrow \, \lnot(\forceQ_1 \Vdash x \in A_{m_{n+1}}))
\end{align*}
but there is a forcing $\forceQ_2$ such that $W[G_{\beta}]$ thinks that 
\begin{align*}
\forceQ_{2} \text{ is } n \text{-legal with respect to $E$ and } \\   {\forceQ_2} \Vdash x \in A_{k_{n+1}}
\end{align*}


Then continuing to argue in $W[G_{\beta}]$, we force with
\[\forceP(\beta):= \hbox{Code}((x,0,m_{n+1},k_{n+1}),2).\] 




\item If neither 1 nor 2 is true, then either \[ \forceP(\beta)=\hbox{Code}((x,0,m_{n+1},k_{n+1}),2)\] or 
\[ \forceP(\beta)=\hbox{Code}((x,0,m_{n+1},k_{n+1}),1)\] depending on whether $ i\in\{1,2\}$ in $F(\beta)$ was 1 or 2. Otherwise $\forceP$ uses the trivial forcing at that stage. 

\item 
If $F(\beta) = (\dot{x},m,k,i)$ and for every $\forceP_{\beta}$-generic filter $G$, $W[G] \models \forall l \le n+1 ((m_l,k_l) \ne (m,k))$,
then let
\[ \forceP(\beta)=\hbox{Code}((x,0,m,k),i)\] depending on whether $ i\in\{1,2\}$ in $F(\beta)$ was 1 or 2.  
\end{enumerate} 
This ends the definition for the successor step $n \rightarrow n+1$.

With this new notion of $n$-legality, we start the proof of the theorem. The ground model over which we form an iteration is the universe $W$ again, which was defined earlier. Over $W$ we will perform first an $\omega$-length, finitely supported iteration $(\forceP_n)_{n \in \omega}$ of legal posets, and then a second legal iteration of length $\omega_1$. The codes of the form $(x,0,m,k)$ shall eventually define the separating sets for $\varphi_m$ and $\varphi_k$; codes of the form $(x,1,m,k)$ shall correspond to countable sets of reals  (i.e. reals themselves) which indicate the correctness of certain codes of the form $(x,0,m,k)$ which avoid the coding areas coded by these reals.

We let $\{(\varphi_{m_n}, \varphi_{k_n}) \, : \, n \in \omega \}$ be an enumeration of all pairs of $\Sigma^1_3$-formulas. Assume that $(\varphi_{m_0}, \varphi_{k_0})$ is such that there is no legal forcing $\forceQ$ such that $W[\forceQ] \models \exists z (\varphi_{m_0}(z) \land \varphi_{k_0}(z))$. Repeating the arguments from before, we set $E_0:=\{m_0,k_0) \}$ and define the notion of 1-legal with respect to $E_0$.
As will become clear in a second, every step of the iteration $(\forceP_n \, : \, {n \in\omega})$ will either use a legal forcing or define a new and gradually stronger notion of legality. We let $l_n \in \omega$ denote the degree of legality we have already defined at stage $n \in \omega$ of our iteration and define $l_0$ to be $0$ (where 0-legal should just be legal) and $l_1$ to be 1 for the base case of our induction; likewise $\forceP_0$ is set to be the trivial forcing.

The forcing $\forceP_{\omega}$ is the countably supported iteration of $(\forceP_n \, : \,n \in \omega)$, which we will define inductively.
Assume we are at stage $n  \ge 1 \in \omega$ of the iteration and we have defined already the following list of objects and notions:
\begin{enumerate}
\item $\forceP_{n-1}$ and the generic filter $G_{n-1}$.
\item A natural number $l_n \le n$ and a notion of $l_n$-legal relative to $E_{l_n}= \{ (m'_0,k'_0),...,(m'_{l_n-1},k'_{l_n-1}) \} \subset \{ (m_0,k_0),...,(m_{n-1},k_{n-1}) \} $, which is a strengthening of 1-legal relative to $E_0$.
\item A finite set of reals  $\{R_0<...<R_{n - (l_n -1)} \}$, where each real $R_i$ codes a countable set of reals. The choice of the indices will become clear later.

\end{enumerate}
Consider now the $n+1$-th pair $(\varphi_{m_n}, \varphi_{k_n})$ and split into cases.

\begin{enumerate}
\item[$(a)$] There is an $l_n$-legal forcing $\forceQ$ such that \[W[G_{n-1}] \models \forceQ \Vdash \exists z (\varphi_{m_n} (z) \land \varphi_{k_n} (z)), \] then we must use the forcing $\forceQ$. We collect all the reals we have added so far generically which witness $({\ast} {\ast} {\ast})$ for a triple $(x,m,k)$ and call the set $R_{n-l_n}$. In a second step, we use the usual method to code the quadruple $(R_{{n-l_n}}, 1,m'_{l_{n-1}},k'_{l_{n-1}} )$ into $\vec{S}^1$. 
\item[$(b)$] In the second case there is no $l_{n}$-legal forcing $\forceQ$ relative to $E_{l_n}$ which forces $\varphi_{m_n}$ and $\varphi_{k_n}$ to have non-empty intersection. In that case we force with the trivial forcing and define the notion $l_{n+1}$-legal. We first let $(m'_{l_n},k'_{l_n})=(m_n,k_n)$ and $E_{l_n +1}:= E_{l_n} \cup \{(m_n,k_n)\}$, and define $l_{n+1}$-legal relative to $E_{l_n+1}$ just as above. 
We  do not define a new $R_{n-l_n}$.
\end{enumerate}

We let $\forceP_{\omega}$ be the inverse limit of the forcings $\forceP_n$ and consider the universe $W[\forceP_{\omega}]$.
We shall and will assume from now on that in $\forceP_{\omega}$, case $(b)$ is applied infinitely many times. 
\begin{lemma}
For every $n \in \omega$, the tail of the iteration $(\forceP_m \, :\, m \ge n)$ is at least an $l_n-1$-legal iteration relative to $E_{l_n}$ as seen from $W[G_n]$.
\end{lemma}
\begin{proof}
This can easily  be seen if one stares at the definition of the iteration. 
At stage $n$ we either follow case $(a)$ which is an $l_n -1$-legal forcing, as we use the $l_n$-legal $\forceQ$ and additionally code $(R_{{n-l_n}}, 1,m'_{l_{n-1}},k'_{l_{n-1}} )$ which results in an $l_n-1$-legal forcing. Or we apply case $(b)$, thus define $l_n+1$-legality and every further factor of the iteration must be $l_n$-legal. As mixed support iterations of $l_{n}-1$-legal forcings yield an $l_n -1$-legal forcing, this ends the proof.

\end{proof}

The arguments which are about to follow will depend in detail heavily on the actual form of the iteration $\forceP_{\omega}$ which in turn depends on how the enumeration of the $\Sigma^1_3$-formulas behaves. The theorem will of course be true, no matter how $\forceP_{\omega}$ does look like. To facilitate the arguments and notation, however, we assume without loss of generality from now on that the sequence $(\varphi_{m_n}, \varphi_{k_n})$ is so chosen that in the definition of $\forceP_{\omega}$ the cases $(a)$ and $(b)$ are alternating. Thus whenever $n$ is even then $(\varphi_{m_n}, \varphi_{k_n})$ is such that case $(b)$ has to be applied and whenever $n$ is odd then for $(\varphi_{m_n}, \varphi_{k_n})$ case $(a)$ is the one which one should apply. Via changing the order of $(\varphi_{m_n}, \varphi_{k_n})$, this can always be achieved. Consequentially, at even stages $2n$ of the iteration, we define the new notion of $n+1$-legal, while forcing with the trivial forcing; on odd stages $2n+1$, we use the $n+1$-legal forcing to force that $\varphi_{m_{2n+1}}$ and $\varphi_{k_{2n+1}}$ have non-empty intersection and then form the real $R_{2n+1 -n} =R_{n+1}$ and then  code the quadruple $(R_{n+1},1,m_{2n},k_{2n})$ into $\vec{S^1}$.

A consequence of the last lemma (and our assumption on the form of $\forceP_{\omega}$) is that for every even natural number $2n$, there is a final stage in $\forceP_{\omega}$ where we create new codes of the form $(R,1,m_{n},k_{n})$. Indeed, by the definition of $l_{n+1}$-legal, no codes of the form $(R,1,m_{n},k_{n})$ are added by $l_{n+1}$-legal forcings and we have that
\begin{align*}
R_{n}:=  \{ R \, : \,  (R, 1,m_{2n},k_{2n}) \text{ is coded into } \vec{S^1} \}.
\end{align*}
To introduce a useful notion, we say that a real $x$, which is coded into $\vec{S^1}$ or $\vec{S^2}$, has coding area almost disjoint from the real $R$, if $R$ codes an $\omega_1$-sized subset of $\omega$ and the $\omega$-blocks, where $x$ is coded are almost disjoint from the set of ordinals coded by $R$, in that their intersection is countable.

\begin{lemma}
In $W[\forceP_{\omega}]$ for every $n \in \omega \backslash 0$, there is only one real $R$, namely $R_{n}$ which has a code of the form $(R_n,1,m_{2n},k_{2n})$ written into $\omega_1$-many  $\omega$-blocks of elements of $\vec{S^1}$ almost disjoint from $R_{n-1}$.
\end{lemma}
\begin{proof}
This is just a straightforward consequence of the definition of the iteration and of our assumption on the form of $\forceP_{\omega}$.  For $n=1$, note that $R_0$ is the unique ordinal for which $(R,1,m_0,k_0)$ is coded into $\vec{S^1}$. Then the next forcing $\forceP(2)$ is trivial while 2-legal is defined and the forcing $\forceP(3)$ first forces $\varphi_{m_3}$ and $\varphi_{k_3}$ to intersect without creating any code of the form $(R,1,m_0,k_0)$ or $(R,1,m_2,k_2)$, forms $R_1$ and writes $(R_1,1,m_2,k_2)$ into $\vec{S}^1$ with coding area almost disjoint from $R_0$. As all later factors of the iteration are 2-legal, there will be no new codes of the form $(R,1,m_2,k_2)$, hence $R_1$ is the unique real for which $(R_1,1,m_2,k_2)$ is written into $\vec{S^1}$ withwith coding area almost disjoint from $R_0$.

The argument for arbitrary $n$ works exactly the same way with the obvious replacements of letters.

\end{proof}
\begin{lemma}
In $W[\forceP_{\omega}]$, every real  $R_{n}$ is $\Sigma^1_3$-definable.
\end{lemma}
\begin{proof}
This is by induction on $n$. For $n=0$, $R_0$ is the unique real  for which $(R,1,m_0,k_0)$ is coded into $\vec{S^1}$.
This can be written in a $\Sigma^1_3$-way:
\begin{align*}
x= R_0 \Leftrightarrow \exists r \forall M &(|M|=\aleph_0 \land M \text{ is transitive } \land r \in M \land \omega_1^M=(\omega_1^L)^M \rightarrow \\&M \text{ sees with the help of the coded information in } r \text{ that }\\& x \text{ is the unique real such that } (x,1,m_0,k_0) \text{ is coded in} \\& \text{ a block of } \vec{S^1}).
\end{align*}

Now assume that there is a $\Sigma ^1_3$-formula which uniquely defines $R_{n}$, then, by the last lemma, $R_{n+1}$ is the unique real which has a code of the form $(R,1,m_{2(n+1)},k_{2(n+1)})$ written into $\aleph_1$-many $\omega$-blocks of elements of $\vec{S^1}$ almost disjoint from the set of ordinals coded by $R_n$. Let $\psi$ be the $\Sigma^1_3$-formula which defines $R_{n}$, then
\begin{align*}
x= \alpha_{n+1} \Leftrightarrow  &\psi(R_{n}) \text{ and } \\ &\exists r \forall M (|M|=\aleph_0 \land M \text{ is transitive } \land r \in M \land \omega_1^M=(\omega_1^L)^M \rightarrow \\&M \text{ sees with the help of the coded information in } r \text{ that }\\& x \text{ is the unique real such that } (x,1,m_{2(n+1)},k_{2(n+1)}) \\& \text{is coded in} \text{ $\aleph_1$-many blocks of } \vec{S^1} \\& \text{almost disjoint from the ordinals coded by }  R_{n}).
\end{align*}

\end{proof}
The ordinals $R_{n}$ indicate the set of places we need to exclude in order ro obtain correct codes of the form $(x,0,m_{2n},k_{2n})$ which are written into $\vec{S}$. Indeed the iteration $\forceP_{\omega}$, after the stage where we coded the quadruple $(R_{n},1,m_{2n},k_{2n})$ into $\vec{S^1}$ will be $2n-1$-legal, just by the definition of the iteration, which means that the tail of $\forceP_{\omega}$ will never produce a pathological situation.

Finally we can form our desired universe of the $\Sigma^1_3$-separation property. We let $E_{\omega} = \{(m_{2n},k_{2n}) \,: \, n \in \omega\}$ and force with $W[\forceP_{\omega}]$ as our ground model.  We use a countably supported iteration of length $\omega_1$ where we force every quadruple of the form $(x,0,m_{2n},k_{2n})$, $x \in 2^{\omega}$, into either $\vec{S^1}$ or $\vec{S^2}$ with coding area almost disjoint from $ R_{n}$, according to whether case 2, 3 or 4 is true in the definition of $n$-legal. Note that, by assumption, all pairs $(\varphi_{m_n}, \varphi_{k_n})$, $n$ odd, do have a non-empty intersection.  Let $W_1$ denote the universe we obtain this way.

As a consequence we can define in $W_1$ the desired separating sets as follows.
For a pair $(\varphi_{m_{2n}},\varphi_{k_{2n}})$ we let $\psi_{n}$ be the $\Sigma^1_3$-formula which defines $\alpha_{n}$. Now for any real $x$  we let 
\begin{align*}
x \in D_{m_{2n},k_{2n}} \Leftrightarrow \, &\psi(R_{n}) \text{ and } \\& \exists r \forall M (|M|=\aleph_0 \land M \text{ is transitive } \land r \in M \land \omega_1^M=(\omega_1^L)^M  \\&\rightarrow M \text{ sees with the help of the coded information in } r \text{ that }\\& (x,0,m_{2n},k_{2n})  \text{ is coded into $\aleph_1$-many blocks of } \vec{S^1} \\& \text{almost disjoint from the ordinals coded by }  R_{n}).
\end{align*}
and
\begin{align*}
x \notin D_{m_{2n},k_{2n}} \Leftrightarrow \, &\psi(R_{n}) \text{ and } \\& \exists r \forall M (|M|=\aleph_0 \land M \text{ is transitive } \land r \in M \land \omega_1^M=(\omega_1^L)^M  \\&\rightarrow M \text{ sees with the help of the coded information in } r \\& \text{that } (x,0,m_{2n},k_{2n})  \text{ is coded into $\aleph_1$-many blocks of } \vec{S^2} \\&\text{almost disjoint from the ordinals coded by }  R_{n}).
\end{align*}
Both formulas are $\Sigma^1_3$, hence the $\Sigma^1_3$-separation property holds in $W_1$.

\end{proof}


\begin{thebibliography}{99}
\bibitem{FS}
S. D. Friedman and D. Schrittesser \textit{Projective Measure without Projective Baire}. Memoirs of the AMS 1298, 2020.
\end{thebibliography}
\end{document}
